至于 (2.2.2), 对图 (2.2.5) 取上同调, 得交换图
\begin{equation}\tag{2.2.19}
\begin{tikzcd}[row sep=huge,column sep=huge]
H^p(\widetilde X,\mathbb Z_\ell)
  \arrow[r,"i^*"]
&
H^p(\widetilde\Delta,\mathbb Z_\ell)
\\
H^p(X,\mathbb Z_\ell)
  \arrow[r,"j^*"]
  \arrow[u,"f^*"']
&
H^p(\Delta,\mathbb Z_\ell)
  \arrow[u,"g^*"']
\end{tikzcd},
\end{equation}
从而得到态射
\begin{equation}\tag{2.2.20}
\frac{H^p(\widetilde X,\mathbb Z_\ell)}
     {f^*H^p(X,\mathbb Z_\ell)}
\xrightarrow{\ (i^*,j^*)\ }
\frac{H^p(\widetilde\Delta,\mathbb Z_\ell)}
     {g^*H^p(\Delta,\mathbb Z_\ell)}.
\end{equation}
以 $F^i$ 表示 $H^\bullet(\widetilde X,\mathbb Z_\ell)$ 和 $H^\bullet(\widetilde\Delta,\mathbb Z_\ell)$ 上由 Leray 谱序列给出的滤过. 则 (2.2.20) 可改写为\footnote{在 (2.2.21) 右侧的分母中, 印本将记号颠倒为 $H^pF^p$. 前一句和左侧所用的滤过是 $F^pH^p$; 此处恢复该记号.}
\begin{equation}\tag{2.2.21}
\frac{H^p(\widetilde X,\mathbb Z_\ell)}
     {F^pH^p(\widetilde X,\mathbb Z_\ell)}
\xrightarrow{\ (i^*,j^*)\ }
\frac{H^p(\widetilde\Delta,\mathbb Z_\ell)}
     {F^pH^p(\widetilde\Delta,\mathbb Z_\ell)}.
\end{equation}
由 (2.2.9), (2.2.21) 在相伴分次对象上诱导的映射为同构, 所以 (2.2.20) 为同构. 另一方面, (2.2.18) 给出直和分解
\begin{equation}\tag{2.2.22}
\begin{tikzcd}[column sep=7em]
H^\bullet(\widetilde X,\mathbb Z_\ell)
  \arrow[r,shift right=1.2ex,
    "f_*\oplus(1-f^*f_*)"']
&
H^\bullet(X,\mathbb Z_\ell)\oplus\Ker f_*
  \arrow[l,shift right=1.2ex,
    "f^*+\text{包含映射}"']
\end{tikzcd},
\end{equation}
而 (1.2.2) 给出直和分解
\begin{equation}\tag{2.2.23}
\frac{H^\bullet(\widetilde\Delta,\mathbb Z_\ell)}
     {g^*H^\bullet(\Delta,\mathbb Z_\ell)}
\xleftarrow{\ \sim\ }
\bigoplus_{1\leq j\leq r}
H^{\bullet-2j}(\Delta,\mathbb Z_\ell(-j)).
\end{equation}
将同构 (2.2.20) 与分解 (2.2.22)、(2.2.23) 结合, 就得到所需的 (2.2.2).

\subsubsection*{注 2.2.24.}
我们略述 (2.2.2) 的第二个证明, 这个证明明显更简单, 适用于 $X$ 在 $k$ 上\underline{固有}的情形. 更一般地, 无须对 $X$ 作任何假设, 在固有支集上同调中可得到相应的结果 (2.2.26). 对 $(\widetilde X,\widetilde\Delta)$ 和 $(X,\Delta)$ 的固有支集上同调 (记为 $H_c^\bullet$) 正合列构成交换图
\begin{center}
\resizebox{\textwidth}{!}{%
\begin{tikzcd}[ampersand replacement=\&,column sep=small,row sep=large]
(2.2.25) \& {} \arrow[r] \&
H_c^i(\widetilde X-\widetilde\Delta,\mathbb Z_\ell)
  \arrow[r] \&
H_c^i(\widetilde X,\mathbb Z_\ell)
  \arrow[r] \&
H_c^i(\widetilde\Delta,\mathbb Z_\ell)
  \arrow[r] \&
H_c^{i+1}(\widetilde X-\widetilde\Delta,\mathbb Z_\ell)
\\
(2.2.26) \& {} \arrow[r] \&
H_c^i(X-\Delta,\mathbb Z_\ell)
  \arrow[r]
  \arrow[u,"\wr" left] \&
H_c^i(X,\mathbb Z_\ell)
  \arrow[r]
  \arrow[u,hook,"f^*" right] \&
H_c^i(\Delta,\mathbb Z_\ell)
  \arrow[r]
  \arrow[u,hook,"g^*" right] \&
H_c^{i+1}(X-\Delta,\mathbb Z_\ell)
  \arrow[u,"\wr" left]
\end{tikzcd}}
\end{center}
无上同调复形 (2.2.25) 对其无上同调子复形 (2.2.26) 的商给出\underline{无上同调}复形
\begin{equation}\tag{2.2.27}
\longrightarrow 0\longrightarrow
\frac{H_c^i(\widetilde X,\mathbb Z_\ell)}
     {f^*H_c^i(X,\mathbb Z_\ell)}
\longrightarrow
\frac{H_c^i(\widetilde\Delta,\mathbb Z_\ell)}
     {g^*H_c^i(\Delta,\mathbb Z_\ell)}
\longrightarrow0\longrightarrow.
\end{equation}
因此, 无须考虑谱序列, 即知 (2.2.20) 为同构. 随后应用刚才的论证, 便得到 (2.2.2).

\subsubsection*{推论 2.3.}
\underline{置}
\[
b_q(X)=\dim_{\mathbb Q_\ell}H^q(X,\mathbb Q_\ell),
\qquad
\chi(X)=\sum_q(-1)^qb_q(X).
\]
\underline{则}\footnote{在 (2.3.1) 中, 印本写作 $b_{q-j}(\Delta)$. 分解 (2.2.2) 将上同调次数移动 $2j$, 所以正确的直和项为 $b_{q-2j}(\Delta)$.}
\begin{equation}\tag{2.3.1}
b_q(\widetilde X)
=b_q(X)+\sum_{j=1}^{r}b_{q-2j}(\Delta),
\end{equation}
\begin{equation}\tag{2.3.2}
\chi(\widetilde X)=\chi(X)+r\,\chi(\Delta).
\end{equation}
